\section{Aggregate-coupling calculations}
\label{app:technical-aggregate}
This section supplies the matrix arithmetic for Proposition~\ref{prop:aggregate_local}
above. For $Q=qI+\gamma\mathbf1\mathbf1^\top$,
\[
Q^{-1}=\frac1q\left(I-\frac{\gamma}{q+\gamma n}\mathbf1\mathbf1^\top\right).
\]
\begin{proof}
Recall the coupled normal equations \eqref{eq:block_conditional_normal_equations}:
\[
 Q_{ii}u_i+\sum_{j\ne i}Q_{ij}\mathbb E[u_j\mid\mathcal G_i]
 =\mathbb E[b_i(\theta_t)\mid\mathcal G_i].
\]
Here $Q_{ii}=q+\gamma$, $Q_{ij}=\gamma$ for $j\ne i$, and $b_i(\theta)=\theta_i$.
Let $\mu_j=\mathbb E[u_j]$. For $j\ne i$, the policy $u_j$ is a function of $\theta_{j,t}$,
which is independent of $\theta_{i,t}$, so
\[
\mathbb E[u_j\mid\mathcal G^{\mathrm{loc}}_{i,t}]
=\mathbb E[u_j]=\mu_j,
\]
while $\theta_{i,t}$ is
$\mathcal G^{\mathrm{loc}}_{i,t}$-measurable. The conditional normal equation
\eqref{eq:block_conditional_normal_equations} therefore gives
\[
(q+\gamma)u_i
+\gamma\sum_{j\ne i}\mu_j
-\theta_{i,t}=0.
\]
Taking expectations yields
\[
q\mu_i+\gamma\sum_j\mu_j=0.
\]
Summing over $i$ implies $\sum_j\mu_j=0$, hence every $\mu_i=0$ and
$u_i=\theta_{i,t}/(q+\gamma)$, which proves
\eqref{eq:optimal_local_shared}. This argument uses independence, zero means,
square integrability, and strict convexity; Gaussianity is not required for the
local-policy formula.

\emph{Open-loop regret.} Here $B=I$ and $\Sigma_\infty=\sigma^2I$, so by \eqref{eq:OU_Rinf}, $R_\infty=\tfrac{\sigma^2}2\operatorname{tr}(Q^{-1})$. From the Sherman--Morrison formula above,

\[
\operatorname{tr}(Q^{-1})=\frac1q\Big(n-\frac{\gamma n}{q+\gamma n}\Big)=\frac{n\,(q+\gamma(n-1))}{q\,(q+\gamma n)},
\]

which gives \eqref{eq:aggregate_Rinf}.

\emph{Local regret.} By the Pythagorean form of Lemma~\ref{lem:radner_projection} (the orthogonality identity after Lemma~\ref{lem:radner_projection}, with both $x^\star$ and $u^{\rm loc}$ mean zero), $2R_{\rm loc}=\|x^\star\|_Q^2-\|u^{\rm loc}\|_Q^2=2R_\infty-\|u^{\rm loc}\|_Q^2$. Since $u^{\rm loc}=\theta/(q+\gamma)$ and $\operatorname{tr}(Q)=n(q+\gamma)$,

\[
\|u^{\rm loc}\|_Q^2=\frac{\sigma^2\operatorname{tr}(Q)}{(q+\gamma)^2}=\frac{n\sigma^2}{q+\gamma}.
\]

Therefore

\[
R_{\rm loc}=\frac{n\sigma^2}{2}\left[\frac{q+\gamma(n-1)}{q(q+\gamma n)}-\frac{1}{q+\gamma}\right]=\frac{n\sigma^2}{2}\cdot\frac{(q+\gamma(n-1))(q+\gamma)-q(q+\gamma n)}{q(q+\gamma)(q+\gamma n)} .
\]

The numerator reduces to $\gamma^2(n-1)$, which gives \eqref{eq:aggregate_Rloc}.

\emph{Share.} Dividing \eqref{eq:aggregate_Rloc} by \eqref{eq:aggregate_Rinf}, the factors $n\sigma^2/2$ and $q(q+\gamma n)$ cancel, leaving \eqref{eq:aggregate_eta}.

\end{proof}

\section{Canonical spatial omission calculation}
\label{app:technical-canonical}
We expand the proof of Lemma~\ref{prop:canonical_eta_S} above,
including the equal-inverse-length case in the variance integral.
\begin{proof}
By spatial stationarity, take $z=0$ and suppress the time argument. Write $\beta_c\triangleq1/\ell_c$ and $\beta_s\triangleq1/\ell_s$ for the inverse lengths, so $w(u)=\tfrac{\beta_c}{2}e^{-\beta_c|u|}$ and $\operatorname{Cov}(\theta(u),\theta(u'))=\sigma^2e^{-\beta_s|u-u'|}$.

\textbf{Step 0 (omission is a variance ratio).} With $Q=qI$ the optimal policy given $\mathcal O_r$ is $\mathbb E[x^\star(0)\mid\mathcal O_r]$. Its expected loss is $\tfrac q2\,\mathbb E\big[\operatorname{Var}(x^\star(0)\mid\mathcal O_r)\big]$. For a Gaussian field the conditional variance is a deterministic number: it does not depend on the observed values (companion Section~\ref{app:toolkit-3}). Since $R_\infty=\tfrac q2\operatorname{Var}(x^\star(0))$,

\[
\eta_S(r)=\frac{\operatorname{Var}(x^\star(0)\mid\mathcal O_r)}{\operatorname{Var}(x^\star(0))}.
\]

\textbf{Step 1 (split into observed middle and two tails).} Write $x^\star(0)=M_r+T_r^++T_r^-$, where

\[
M_r=\int_{-r}^{r}w(u)\theta(u)\,du,\qquad T_r^+=\int_r^\infty w(u)\theta(u)\,du,\qquad T_r^-=\int_{-\infty}^{-r}w(u)\theta(u)\,du .
\]

The middle $M_r$ is computed from observed values, so it contributes no conditional variance.

\textbf{Step 2 (rewrite the right tail).} Substituting $u=r+s$,

\[
T_r^+=\frac{\beta_c}{2}e^{-\beta_c r}\,Y_r,\qquad Y_r\triangleq\int_0^\infty e^{-\beta_c s}\,\theta(r+s)\,ds .
\]

\textbf{Step 3 (spatial Markov property).} By Section~\ref{sec:spatiotemporal_predictability}, the field is Markov in $z$. Given the observed segment $[-r,r]$, the right tail $\{\theta(u):u>r\}$ and the left tail $\{\theta(u):u<-r\}$ are conditionally independent, and each depends on the segment only through its adjacent endpoint $\theta(r)$ or $\theta(-r)$. Therefore

\[
\operatorname{Var}(T_r^+\mid\mathcal O_r)=\operatorname{Var}(T_r^+\mid\theta(r)),\qquad \operatorname{Var}(T_r^-\mid\mathcal O_r)=\operatorname{Var}(T_r^-\mid\theta(-r)),\qquad \operatorname{Cov}(T_r^+,T_r^-\mid\mathcal O_r)=0 .
\]

\textbf{Step 4 (the tail variance does not depend on $r$).} Define $C\triangleq\operatorname{Var}(Y_0\mid\theta(0))$. Shifting the field by $r$ does not change its law, so $\operatorname{Var}(Y_r\mid\theta(r))=C$ for every $r$. By reflection symmetry (the covariance depends only on $|u-u'|$), the left tail gives the same value. Hence

\[
\operatorname{Var}(x^\star(0)\mid\mathcal O_r)=2\Big(\frac{\beta_c}{2}\Big)^2e^{-2\beta_c r}\,C=\frac{\beta_c^2}{2}e^{-2\beta_c r}\,C .
\]

At $r=0$ the same argument applies with $\mathcal O_0=\sigma(\theta(0))$ and $M_0=0$, giving $\operatorname{Var}(x^\star(0)\mid\theta(0))=\tfrac{\beta_c^2}{2}C$. Dividing, $\eta_S(r)=\eta_S(0)\,e^{-2\beta_c r}=\eta_S(0)\,e^{-2r/\ell_c}$. The unknown constant $C$ cancels, so it never has to be computed.

\textbf{Step 5 (the constant $\eta_S(0)$).} Conditioning on the single Gaussian variable $\theta(0)$ gives $\operatorname{Var}(x^\star(0)\mid\theta(0))=\operatorname{Var}(x^\star(0))\big(1-\operatorname{Corr}(x^\star(0),\theta(0))^2\big)$ (companion Section~\ref{app:toolkit-3}), so $\eta_S(0)=1-\operatorname{Corr}^2$. Two moments are needed.

\emph{Covariance with the local value:}

\[
\operatorname{Cov}(x^\star(0),\theta(0))=\int_{\mathbb R}w(u)\,\sigma^2e^{-\beta_s|u|}\,du=\frac{\sigma^2\beta_c}{2}\int_{\mathbb R}e^{-(\beta_c+\beta_s)|u|}\,du=\frac{\sigma^2\beta_c}{\beta_c+\beta_s}.
\]

\emph{Variance of the target:} $\operatorname{Var}(x^\star(0))=\frac{\sigma^2\beta_c^2}{4}\,I$, where $I\triangleq\iint e^{-\beta_c(|u|+|v|)-\beta_s|u-v|}\,du\,dv$.

To evaluate $I$, define the inner integral $h(u)\triangleq\int_{\mathbb R}e^{-\beta_c|v|}e^{-\beta_s|u-v|}\,dv$. Then $h$ is even, so $I=2\int_0^\infty e^{-\beta_c u}h(u)\,du$. For $u\ge0$, split the $v$-integral at $0$ and $u$:

\[
h(u)=\underbrace{\frac{e^{-\beta_s u}}{\beta_c+\beta_s}}_{v<0}+\underbrace{\frac{e^{-\beta_c u}-e^{-\beta_s u}}{\beta_s-\beta_c}}_{0\le v\le u}+\underbrace{\frac{e^{-\beta_c u}}{\beta_c+\beta_s}}_{v>u}.
\]

(The middle term is read as $u\,e^{-\beta_c u}$ when $\beta_s=\beta_c$.) Multiply by $e^{-\beta_c u}$ and integrate over $u\ge0$ term by term:

\[
\frac{1}{(\beta_c+\beta_s)^2},\qquad \frac{1}{\beta_s-\beta_c}\Big(\frac{1}{2\beta_c}-\frac{1}{\beta_c+\beta_s}\Big)=\frac{1}{2\beta_c(\beta_c+\beta_s)},\qquad \frac{1}{2\beta_c(\beta_c+\beta_s)} .
\]

Summing and doubling,

\begin{equation}
I=2\left[\frac{1}{(\beta_c+\beta_s)^2}+\frac{1}{\beta_c(\beta_c+\beta_s)}\right]=\frac{2(2\beta_c+\beta_s)}{\beta_c(\beta_c+\beta_s)^2},\qquad \operatorname{Var}(x^\star(0))=\frac{\sigma^2\beta_c(2\beta_c+\beta_s)}{2(\beta_c+\beta_s)^2}.
\end{equation}

\emph{Combine.} With $\operatorname{Var}(\theta(0))=\sigma^2$,

\[
\operatorname{Corr}^2=\frac{\big(\sigma^2\beta_c/(\beta_c+\beta_s)\big)^2}{\sigma^2\cdot\sigma^2\beta_c(2\beta_c+\beta_s)/\big(2(\beta_c+\beta_s)^2\big)}=\frac{2\beta_c}{2\beta_c+\beta_s},\qquad \eta_S(0)=\frac{\beta_s}{2\beta_c+\beta_s}=\frac{\ell_c}{\ell_c+2\ell_s}.
\]
\end{proof}

\section{Numerical reproducibility details}
\label{app:technical-numerics}
This section gives the settings and estimators for the three figures in
Section~\ref{sec:numerical-evaluation}. The crossover uses the exact coupled
fresh-local rule; the radius curves use scalar variance composition; and
the ring uses agentwise Gaussian conditioning with $Q=I$.
Normalized scalar and ring losses use analytical open-loop regret
$R_\infty$. The crossover instead uses paired sample means.
Pointwise Monte Carlo means have 95\% normal intervals
$\bar L\pm1.96s_L/\sqrt N$, with sample standard deviation $s_L$.
Conditioning uses linear solves rather than explicit matrix inverses.

\begin{table}[H]
\centering
\caption{Settings for the three numerical illustrations. The shared-resource
model fixes $q=1$; the ring covariance has average marginal variance one.}
\label{tab:numerical-parameters}
\small
\begin{tabular}{llll}
\toprule
Figure & principal sweep & fixed setting/size & samples \\
\midrule
\ref{fig:local-global-phase} & 9 $\gamma/q$ values in $[0,10]$ & $n=5,10,25,100$ & $30{,}000$/setting \\
\ref{fig:radius-regret} & $r/\ell_c\in[0,2]$ (1,001 points) & 3 regimes; $D/\ell_c=2$ & $50{,}000$/point \\
\ref{fig:decision-relevance} & $\ell_c=0.7,2,6$ & 64-node ring; $T=6$ & $12{,}000$/check \\
\bottomrule
\end{tabular}
\end{table}

\paragraph{Fresh-local/stale-global crossover (Fig.~\ref{fig:local-global-phase}).}
We fix $q=1$, use $n\in\{5,10,25,100\}$, and sample independent
standard-normal state components ($\sigma^2=1$).
The nine coupling ratios are $\gamma/q\in\{0,0.05,0.1,0.2,0.5,1,2,5,10\}$.
At each setting, 30,000 independent current-state draws give paired open-loop and fresh-local
losses, using $u_i=\theta_i/(q+\gamma)$.
Writing their means as $o=\widehat R_\infty$ and
$l=\widehat R_{\rm loc}$, the estimator is
$\widehat\rho^\star=\tfrac12\log[o/(o-l)]$.
Its 95\% delta-method interval is
$\widehat\rho^\star\pm1.96\sqrt{g^\top Cg/N}$, where $C$ is the sample
covariance of the paired losses and
\[
g=\left(-\frac{l}{2o(o-l)},\,\frac{1}{2(o-l)}\right)^\top.
\]
The paired covariance is retained; $o-l\le0$ is rejected rather than clipped.
The zero-coupling boundary is zero and is omitted from the logarithmic plot.

\paragraph{Three radius regimes (Fig.~\ref{fig:radius-regret}).}
The scalar calculation uses $\eta_S(r)=\eta_0e^{-2r/\ell_c}$,
$\tau(r)=r/v$, $\ell_c=1$, and $D=2$, with
$(\eta_0,vT)=(0.25,0.2),(0.7,4),(0.9,100)$.
The radius grid has spacing 0.002. At each point, 50,000 pairs of independent
standard normals are scaled by the square roots of the normalized components
$R_T/R_\infty$ and $R_S/R_\infty$ in
\eqref{eq:temporal_component}--\eqref{eq:spatial_component}, then added
and squared. This checks scalar variance composition; it does not simulate
a spatial field.

\paragraph{Decision relevance on a ring (Fig.~\ref{fig:decision-relevance}).}
The graph and covariance are fixed on a 64-node ring, with $T=6$ and latency
0.12 per hop. For $\ell_c\in\{0.7,2,6\}$ the operator is
$K_{ij}=e^{-d_G(i,j)/\ell_c}/\|e^{-d_G(i,\cdot)/\ell_c}\|_2$.
With $Q=I$, exact Gaussian conditioning gives spatial omission; 12,000-sample
checks evaluate selected radii. Each operator has its own analytical
$R_\infty=\tfrac12\operatorname{tr}(K\Sigma_S K^\top)$.
Raw state correlation is computed separately from the fixed covariance.

The covariance in Eq.~\eqref{eq:graph_spectral_covariance} uses
$(\kappa_s,\nu)=(0.75,1.5)$ and is scaled to average marginal variance one.
Every integer radius from 0 through the ring diameter 32 is evaluated.
For each operator, Monte Carlo checks use radii
$\{0,r^\star,\min(r^\star+3,32),32\}$. Independent current-state Gaussian
draws are passed through the exact conditional policies; summed squared
prediction errors are divided by $\operatorname{tr}(K\Sigma_SK^\top)$
to estimate the omission share.

\paragraph{Reproduction.}
The \href{https://github.com/ANRGUSC/TINA}{public TINA repository} retains
the original experiment identifiers. Figures~\ref{fig:local-global-phase},
\ref{fig:radius-regret}, and \ref{fig:decision-relevance} correspond to
\path{numerics/config/experiment_02.yaml},
\path{numerics/config/experiment_03.yaml}, and
\path{numerics/config/experiment_05.yaml}, with seeds 20260824, 20260825,
and 20260827, respectively. Their drivers are
\path{numerics/experiments/exp02_local_global_crossover.py},
\path{numerics/experiments/exp03_radius_regret.py}, and
\path{numerics/experiments/exp05_decision_relevance.py}.
Run metadata, processed results, and selected raw arrays are stored under
\path{numerics/results/}; not every random draw is archived.
\path{docs/numerical_method.md} and \path{docs/experiment_parameters.md}
specify conditioning safeguards and environment versions;
\path{numerics/README.md} gives reproduction commands.
